\section{课程导入：多 Agent 不是新概念，但进入了新阶段}

Noam Brown 在开场先把范围说清楚：这讲不是只讲一个特定算法，而是把 \textbf{自博弈、博弈均衡、人类协作、LLM 推理扩展}放到同一个框架里讨论。他强调当前很多团队把多 Agent 当成 ``工程编排技巧''，但从历史上看，它首先是一个 \textbf{决策理论问题}，其次才是系统工程问题。

\begin{importantbox}{本讲主问题}
如果目标是构建现实环境中可部署的 Agent，那么我们到底应该优化什么：Minimax Equilibrium，还是 Population Best Response？
\end{importantbox}

他给出的叙事很有针对性：过去十年游戏 AI 的突破，证明了 ``递归式自改进'' 在特定假设下非常强；但一旦从二人零和、完美信息跳到多人、非零和、带语言协商的环境，原先的漂亮性质会快速失效。

\begin{knowledgebox}{为什么这讲对 LLM Agent 特别关键}
\begin{itemize}
    \item 它把 ``推理时扩展（test-time scaling）'' 和 ``多体交互（multi-agent interaction）'' 放在统一坐标系。
    \item 它解释了为什么一些在 benchmark 上有效的 agent scaffold，放到真实协作场景会失灵。
    \item 它给出一个强烈但可检验的命题：在非二人零和任务中，\textbf{没有目标群体数据，就学不到稳定协作策略}。
\end{itemize}
\end{knowledgebox}

\subsection{本章小结}
这讲不是 ``多 Agent 技巧汇总''，而是一次目标函数审计。先问 ``什么是好策略''，再谈如何训练、如何部署。

